% Completed proof, promoted after targeted root review on 2026-11-20.
\subsection{A generic hyperelliptic Jacobian obstructs integral autoduality}
\paragraph{J2a (Refuted, with a complete counterexample).}
Let $a_0,\ldots,a_6$ be algebraically independent over $\mathbb F_5$,
put $K=\mathbb F_5(a_0,\ldots,a_6)^{\mathrm{perf}}$ and
$k=\overline K$, and let
\[
 C/K:\quad y^2=x^7+a_6x^6+\cdots+a_1x+a_0,
 \qquad A=\operatorname{Jac}(C),\qquad X=A_k.
\]
Here $C$ means the smooth projective model and $\infty$ is its
unique point above infinity.  Its genus is three, and $X$ is a
smooth projective threefold in characteristic five.  We prove
that there is no isomorphism from $\operatorname{Ab}^2(X)$ to its
dual induced by the integral intersection of families of cycles.
We use the precise required compatibility: for every family
$z\in\operatorname{CH}^2(T\times X)$, $z(0)=0$, on an abelian
parameter variety $T$, if $u_z:T\to\operatorname{Ab}^2(X)$ is its
regular homomorphism, the homomorphism defined by
$r_*(p^*z\cdot q^*z)$ must equal $u_z^\vee\rho u_z$.
This is exactly the construction in 44:9; no arbitrary rescaling
of intersection is allowed.

Write $B=\operatorname{Ab}^2(A/K)$.  This abelian variety exists,
and its base change to $k$ is the usual algebraic representative,
by \cite[Theorem 5.10, pp.~25--26; Theorem 6.1, pp.~26--27]{DGJ_ACMVFunctor}.
The compatibility defining $\rho$, if satisfied, determines it
uniquely: applying it to a miniversal family on $B$ gives
$[r]^\vee\rho[r]$ for some $r\ne0$, and $\operatorname{Hom}(B,B^\vee)$
is torsion-free.  Therefore every Galois conjugate of $\rho$ is the
same $\rho$.  It descends to $K$.  The perfectness of $K$ makes
$k/K$ separable, so the stated base-change and descent theorem
applies without an inseparable-field qualification.

\paragraph{The dimension of $B$ is at most three.}
Choose a miniversal family for $B$, and a prime $\ell\ne2,5$
not dividing its degree.  Its rational Tate realization is
injective by \cite[Proposition 8.2, p.~31]{DGJ_ACMVFunctor}.
After dualizing and applying Poincaré duality on the threefold,
this makes $H^1(B,\mathbb Q_\ell)$ a Galois quotient of
$H^3(A,\mathbb Q_\ell)(1)$.
The geometric Zariski monodromy on $V=H^1(A,\mathbb Q_\ell)$ is
$\operatorname{Sp}_6$.  Indeed, this is
\cite[Theorem 3.4 and Corollary 3.5, pp.~13--14]{DGJ_AchterPries2006}
for the hyperelliptic moduli space.  Here is the precise transfer to our coefficient space
$S=\operatorname{Spec}\mathbb F_5[a_0,\ldots,a_6,\Delta^{-1}]$.
Let $\mathcal H_3^w$ be the stack of smooth hyperelliptic
curves with one marked Weierstrass point.  The morphism
$\mathcal H_3^w\to\mathcal H_3$ forgetting the point is finite
étale of degree eight: its pullback to any family is that
family's reduced branch divisor.  Thus its geometric fundamental
group has finite-index image.  Finite-index subgroups of a
Zariski-dense subgroup of the connected group
$\operatorname{Sp}_6$ are still Zariski dense.

There is, more precisely, a presentation
$\mathcal H_3^w=[S/G]$ with connected smooth group
\[
 G=\{(b,a,c):b\in\mathbb G_a,\ a,c\in\mathbb G_m,
                    \ c^2=a^7\}
       \simeq\mathbb G_a\rtimes\mathbb G_m.
\]
The coordinate changes are $x\mapsto ax+b$, $y\mapsto cy$;
the monic leading coefficient imposes $c^2=a^7$.
The isomorphism on the multiplicative factor is
$s\mapsto(a,c)=(s^2,s^7)$, with inverse $s=c/a^3$.
Every pointed hyperelliptic curve has this form fppf locally:
the marked branch point trivializes the quotient conic locally,
we choose it as infinity, and we choose the anti-invariant
coordinate $y$ and normalize its leading coefficient.  Any
isomorphism of two such presentations has exactly the displayed
form, since functions with a pole of order at most two at the
marked point are spanned by $1,x$, and the anti-invariant
functions with pole order at most seven are spanned by $y$.
This proves the claimed presentation, also for families by the
corresponding locally free sheaves and base change.
The atlas $S\to[S/G]$ has geometrically connected fibres $G$,
so it surjects on geometric fundamental groups: a connected
finite étale cover of the stack stays connected after this
pullback, which is the finite-cover criterion for surjectivity.
Consequently the geometric Zariski monodromy on $S$ is
$\operatorname{Sp}_6$.  Removing a proper closed subset of
normal $S$ gives a surjection on fundamental groups, and
purely inseparable extension to $K$ leaves the étale Galois
representations unchanged.

The symplectic exterior-cube decomposition is
\[
 \bigwedge^3V(1)=V\oplus P,
 \qquad \dim V=6,\quad \dim P=14,
\]
where the displayed $V$ denotes its Lefschetz copy and $P$ is
primitive.  Both summands are absolutely irreducible for the
geometric symplectic group and are preserved by the arithmetic
Galois action.  A direct verification of the only less familiar
assertion is as follows.  With symplectic basis $e_i,f_i$, $P$
is the kernel of symplectic contraction.  Its fourteen distinct
torus weights consist of the eight $\pm\epsilon_1\pm\epsilon_2
\pm\epsilon_3$ and the six $\pm\epsilon_i$.  The first weight
vectors are wedges choosing one of $e_i,f_i$ for each $i$; the
second are $v_i\wedge(e_j\wedge f_j-e_h\wedge f_h)$ for
$\{i,j,h\}=\{1,2,3\}$ and $v_i=e_i$ or $f_i$.
Symplectic permutations and interchanges generate all eight
vectors of the first kind from any one.  For an explicit root operation, let $N(e_2)=e_3$,
$N(f_3)=-f_2$, and let $N$ vanish on the remaining basis
vectors.  It belongs to the symplectic Lie algebra and sends
$e_1\wedge e_2\wedge f_3$ to
$-e_1\wedge(e_2\wedge f_2-e_3\wedge f_3)$.
The opposite root $N'(e_3)=e_2$, $N'(f_2)=-f_3$ sends the
latter second-kind vector to $-2e_1\wedge e_2\wedge f_3$.
Permuting and interchanging symplectic pairs gives all six
second-kind vectors and the converse transitions.  Thus a nonzero
subrepresentation, which contains a torus weight vector, contains
all fourteen vectors.  This proves irreducibility over
$\mathbb Q_\ell$, where $2\ne0$.
Consequently the quotient $H^1(B)$ can have only the constituents
$V$ and $P$, each at most once.

The constituent $P$ cannot occur.  We give the spreading
step explicitly.  Write
$K_n=\mathbb F_5(a_0^{1/5^n},\ldots,a_6^{1/5^n})$, so
$K=\bigcup_nK_n$.  The abelian variety $B$, its group laws,
and the finitely many components and coefficients of a
miniversal cycle $\Gamma\in\operatorname{CH}^2(B\times A)$
are of finite presentation; hence they descend to some $K_n$.
The model of $A$ is the Frobenius pullback of its model over
$S$.  After removing a proper closed subset of this Frobenius
pullback $S_n$, the descended $B$ extends to an abelian scheme
$\mathcal B\to U$, and the closures of the cycle components
give a cycle on $\mathcal B\times_U\mathcal A$.  To obtain
the abelian scheme, extend projective equations and the group
laws and their inverse, shrink so that the group identities
hold, and shrink again to the smooth locus with geometrically
connected fibres.  Properness is retained by the projective
model.  No extension of the universal regular property to every
fibre is being assumed.

The generic Tate injection of Proposition~8.2 is equivariant
for $\operatorname{Gal}(K_n)$: purely inseparable base change
identifies this group and its étale cohomology with those over
$K$.  Both its source and target are the generic fibres of
lisse $\mathbb Q_\ell$-sheaves on $U$, and therefore factor
through $\pi_1(U)$.  Since
$\operatorname{Gal}(K_n)\to\pi_1(U)$ is surjective, their
equivariant map is a map of those lisse sheaves.  Its rank is
locally constant.  Dualizing gives a surjection on every fibre
of $U$, just as at the generic point.  Its kernel is generically
one of the sums of the two intrinsic Lefschetz summands;
equality of morphisms of lisse sheaves at a geometric generic
point is equality everywhere, so its fibrewise quotient has the
same summands.  The Frobenius pullback and shrinking preserve
connected geometric monodromy by the preceding fundamental-group
argument.  The ordinary locus of the Jacobian is
nonempty.  For an explicit check, the smooth curve
$y^2=x^7-x$ over $\mathbb F_5$ has Cartier--Manin matrix
\[
 \begin{pmatrix}0&0&1\\0&3&0\\1&0&0\end{pmatrix},
 \qquad\det=2\pmod5.
\]
This follows directly from applying Cartier to
$x^{j-1}dx/y=y^{-5}x^{j-1}(x^7-x)^2dx$:
the $(i,j)$ entry is the coefficient of $x^{5i-j}$ in
$(x^7-x)^2$.  Thus Cartier is invertible, so the Jacobian is
ordinary.  Nonvanishing of the same determinant defines a
nonempty open ordinary locus in the coefficient space.  It
meets the open set where the preceding objects have spread;
choose a closed point there, over a finite field $\mathbb F_q$.

For clarity, the passage from the Cartier calculation to
Frobenius slopes can be made entirely in terms of characteristic
polynomials.  The Manin congruence in
\cite[\S2.5, p.~7; Cartier formula in \S3.1, pp.~7--8]{DGJ_AchterHowe2018}
says that the Frobenius polynomial modulo five is $T^3$ times
the characteristic polynomial of an invertible iterate of the
Cartier operator (equivalently its adjoint Frobenius operator).
That second polynomial has nonzero constant coefficient.
Hensel factorization therefore gives exactly three unit roots.
The principal polarization pairs each root $\alpha$ with
$q/\alpha$; see also \cite[Lemma 19.2, p.~41]{DGJ_MilneAV}.
Hence the remaining three roots have normalized valuation one.
Thus the six Frobenius eigenvalues on $H^1(A)$ have three
$q$-adic slopes zero and three slopes one.  Their triple
products on $H^3(A)(1)$ include one slope $-1$ and one slope $2$.
The Lefschetz summand has only slopes zero and one, so both
extreme slopes belong to $P$.  The surjection on $\ell$-adic fibres gives equality of
Frobenius characteristic polynomials for $H^1(B)$ and the
indicated summand or sum of summands.  These polynomials have
rational coefficients: the primitive polynomial is the quotient
of the exterior-cube polynomial by the Lefschetz polynomial.
Equality in $\mathbb Q_\ell[T]$ is therefore equality in
$\mathbb Q[T]$.  Embed their algebraic roots in
$\overline{\mathbb Q}_5$ and normalize $v(q)=1$.
If $P$ occurred in the quotient $H^1(B)$, its negative-valuation
root would then be a Frobenius eigenvalue on $H^1(B)$.  This
contradicts the algebraic integrality of Frobenius eigenvalues
on abelian varieties, as stated in
\cite[Theorem 19.1 and its proof, pp.~40--41]{DGJ_MilneAV}.
This argument compares characteristic polynomials; it makes no
assertion that the correspondence admits a particular crystalline
realization.  Hence only the
six-dimensional constituent can occur, and $\dim B\leq3$.
The finite computation records the full primitive slope
multiplicities $1,6,6,1$ in slopes $-1,0,1,2$; the exclusion uses
only the first of these.

\paragraph{The translation family has intersection $\pm2\lambda_A$.}
Embed $C$ by $Q\mapsto[Q-\infty]$ and let
\[
 z_t=[C+t]-[C],\qquad t\in A.
\]
Equivalently $z=m_-^*[C]-\operatorname{pr}_X^*[C]$ on
$A\times A$, where $m_-(t,x)=x-t$.  This is a specified family
of algebraically trivial codimension-two cycles with $z_0=0$.
It induces a $K$-homomorphism $u:A\to B$.
Let $\theta$ denote the cohomology class of the principal theta
polarization.  The Poincaré formula is $[C]=\theta^2/2$ in
cohomology.  It follows as well by characterizing $[C]$ by
$\int_A[C]\alpha\beta=\int_Ca^*\alpha\,a^*\beta$ and the
symplectic pairing defining the principal Jacobian polarization.
With symplectic basis $e_i,f_i$ and
$\theta=\sum e_i\wedge f_i$, the part of the class of $z$ in
$H^1(A_{\rm par})\otimes H^3(A_{\rm space})$ is obtained by
expanding
\[
 \frac12\left(\sum_{i=1}^3
 (e_{i,x}-e_{i,t})\wedge(f_{i,x}-f_{i,t})\right)^2
       -\frac12\left(\sum e_{i,x}\wedge f_{i,x}\right)^2.
\]
All terms of parameter degree zero cancel.  Therefore in
$r_*(p^*z\cdot q^*z)$ only the product of the two parameter-degree-one
parts contributes.  Wedge multiplication and integration of
$e_{1,x}f_{1,x}e_{2,x}f_{2,x}e_{3,x}f_{3,x}$ give
\[
 -2\sum_{i=1}^3
 (e_{i,t}\wedge f_{i,s}-f_{i,t}\wedge e_{i,s}).
\]
This is $-2$ times the normalized Poincaré pairing of the
principal polarization, with the sign reversed if the opposite
Poincaré convention is chosen.  Faithfulness of the Tate
realization on homomorphisms of abelian varieties
\cite[Theorem 12.5, pp.~21--22]{DGJ_MilneAV} gives the
actual homomorphism $\pm2\lambda_A$, not merely its isogeny class.
The wedge calculation is also saved as an exact integer
calculation in \texttt{check\_generic\_jacobian.py}.

\paragraph{The impossible degree.}
If an intersection-induced isomorphism $\rho:B_k\to B_k^\vee$
existed, its defining compatibility and the preceding calculation
would give
\[
 u^\vee\rho u=\pm2\lambda_A.
\]
The right side is an isogeny.  Thus $u$ has finite kernel;
since $\dim B\leq3=\dim A$, it is an isogeny onto $B$.
Taking degrees yields
\[
 (\deg u)^2=\deg(2\lambda_A)=2^6,
 \qquad\deg u=8.
\]
We show that the $K$-abelian variety $A$ admits no isogeny of
degree eight.  The splitting field of its generic septic has
Galois group $S_7$: equivalently, start with seven independent
roots and take the invariant field of their permutations.
Passing to the perfect closure is purely inseparable and leaves
this group unchanged.  If $r_1,\ldots,r_7$ are the roots and $W_1,\ldots,W_7$ are the finite
Weierstrass points, the elements $[W_i-\infty]$ generate $A[2](k)$,
with their sum as the only relation.  To check the last assertion,
a putative proper-subset relation would make
$\prod_{i\in I}(x-r_i)$ a square in $k(C)$.  In the quadratic
extension $k(C)/k(x)$, a rational function whose square lies in
$k(x)$ is in $k(x)$ or in $y\,k(x)$.  The odd valuations at the
chosen roots exclude the first case unless $I$ is empty and
exclude the second unless $I$ is the full set.  There are
therefore six independent classes, as required by the rank of
$A[2]$.

It follows that $A[2](k)$ is the six-dimensional $S_7$-module
$\mathbb F_2^7/\langle(1,\ldots,1)\rangle$, equivalently the
subspace of vectors with even coordinate sum.  This module is
irreducible.  Indeed any nonzero even-sum vector has both a zero
and a one coordinate; subtracting its image under the
transposition of those coordinates yields a vector supported
on exactly two positions.  The $S_7$-orbit of that vector spans
the whole even-sum subspace.  Every nonzero invariant subspace
is consequently the whole module.

The kernel of a degree-eight isogeny in characteristic five is
an étale group of order eight.  Its subgroup killed by two is a
nonzero Galois-stable subgroup of $A[2](k)$, and has at most
eight elements.  Irreducibility would force it to be all of
$A[2](k)$, which has sixty-four elements.  This contradiction
proves the claimed impossibility of an integral
intersection-induced isomorphism.

This argument uses only the rational miniversal Tate embedding,
not the unproved integral torsion standard assumption.  It
rules out the raw integral autoduality demanded in J.2.  It does
not rule out a rational pairing, an isogeny, a polarization after
an appropriate Lefschetz normalization, or an unrelated abstract
isomorphism of $B$ with its dual.
